% Generated draft: requires mathematical and source review.
\documentclass[11pt]{article}
\usepackage[margin=25mm]{geometry}
\usepackage{fontspec}
\setmainfont{Latin Modern Roman}
\usepackage{amsmath,amssymb,mathtools,unicode-math}
\setmathfont{Latin Modern Math}
\usepackage{enumitem,needspace,xurl,hyperref}
\setlength{\parindent}{0pt}
\setlength{\parskip}{5pt}
\setlength{\emergencystretch}{3em}
\newcommand{\N}{\mathbb{N}}
\newcommand{\Z}{\mathbb{Z}}
\newcommand{\Q}{\mathbb{Q}}
\newcommand{\R}{\mathbb{R}}
\newcommand{\C}{\mathbb{C}}
\newcommand{\F}{\mathbb{F}}
\newcommand{\A}{\mathbb{A}}
\newcommand{\PP}{\mathbb{P}}
\newcommand{\E}{\mathbb{E}}
\newcommand{\eps}{\varepsilon}
\newcommand{\dd}{\,\mathrm{d}}
\begin{document}
\section*{UnsolvedMath ID: 30001736}
\textbf{Source problem code:} catalogue entry
\subsection{Definitions and mathematical statement}
Let \(\theta\colon G\to G\) be an involutive automorphism, so that \(\theta^2=\operatorname{id}_G\), and let
\[
H=G^\theta:=\{g\in G:\theta(g)=g\}.
\]
Here \(H\) denotes the full fixed-point subgroup specified in the source. For a complex representation \(\pi\) of \(G\) on a nonzero space \(V_\pi\), give \(\C\) the trivial \(H\)-action. An \(H\)-invariant linear functional satisfies
\[
\ell(\pi(h)v)=\ell(v)\qquad(h\in H,\ v\in V_\pi).
\]
The notation \(\operatorname{Hom}_H(\pi,\C)\) denotes the space of such functionals, with any continuity or other regularity requirements imposed by the intended representation category. The representation \(\pi\) is called \(H\)-distinguished when
\[
\operatorname{Hom}_H(\pi,\C)\ne\{0\}.
\]
\textbf{Classification question.} For a symmetric pair \(H=G^\theta\), determine which irreducible representations \(\pi\) of \(G\) are \(H\)-distinguished. Equivalently, once the group setting and representation category have been specified, characterize the isomorphism classes
\[
\bigl\{[\pi]:\pi\text{ is an irreducible representation of }G,
\ \operatorname{Hom}_H(\pi,\C)\ne\{0\}\bigr\}.
\]
\textbf{Unresolved scope.} The supplied record does not specify the class of groups \(G\), an underlying field or topology, or the representation category. It also does not specify whether invariant functionals must be continuous. Consequently, this is a conditional classification schema preserving the source question; it does not supply a fully specified classification problem in any particular setting.
\subsection{Short English statement}
For the full fixed-point subgroup \(H=G^\theta\) of an involution, classify the irreducible representations of \(G\) admitting a nonzero \(H\)-invariant complex linear functional. The group setting, representation category, and required regularity of the functional remain unspecified in the supplied record.
\subsection{Sources}
\begin{itemize}
\item Local UnsolvedMath registry, \texttt{data/unsolvedmath/problems.json}, permanent ID \(30001736\), title ``Distinguished Representations of Symmetric Pairs,'' source code \texttt{catalogue entry}; collection: Oberwolfach Reports Open Problems.
\item Supplied source citation: \emph{Automorphic Forms: New Directions} (2011), Oberwolfach Reports, pp.\ 719--779, DOI 10.4171/owr/2011/14. \url{https://doi.org/10.4171/owr/2011/14}. The record attributes the question to the workshop-report collection and gives proposed year 2011, but names no individual proposer or precise location within the report.
\item Supplied literature source, labeled ``symmetric-variety dual groups'': \url{https://arxiv.org/abs/2308.15800}. No authors, title, or specific supporting passage are supplied for this link.
\item Literature assessment dated exactly \texttt{2026-08-21}: the supplied curation labels the problem open and describes classification for general symmetric pairs as active, with conjectures concerning relations to symmetric-space dual groups. This is supplied dated curation, not an independently verified assessment; the record explicitly requests independent verification of time-sensitive status and cited claims.
\item The registry labels the statement \texttt{corrected\_verified} and reports that its mathematical content was checked against the official Oberwolfach report, with notation, context, and question form normalized without changing the source-level claim. This is a supplied verification claim; no independent source consultation is asserted here.
\end{itemize}
\end{document}
